% ===========================================================================
% Section 8: Discussion  +  Section 9: Conclusion  (W15 + W16)
%
% Discussion is the designated interpretation section: Robot-mode facts
% live in Sections 5-7; here interpretation is allowed but every claim
% must trace to a table, figure, or the replication archive (footnoted).
%
% Four subsections per tasks.md W15:
%   8.1 Fiscal consolidation as regional policy
%   8.2 MA elasticities in comparative perspective (the Gibbons gap,
%       sharpened by the diagnostics: anchor variants, spatial SEs,
%       theta sensitivity)
%   8.3 Sector-mode specialization (B&P scale economies)
%   8.4 Limitations
%
% Numbers verified against: table CSVs (9, 10, 13, 14), Plan memo /
% diagnostics outputs (theta sweep, refpoint variants, Conley).
%
% CITATION KEYS added this PR: gibbons2024 (STUB), ahlfeldtfeddersen2018.
% ===========================================================================

\section{Discussion}
\label{sec:discussion}

\subsection{Fiscal Consolidation as Regional Policy}
\label{sec:fiscal_policy}

The reform studied here was not designed as regional policy. Section~2
documents that the closures followed fiscal criteria---operating
deficits and traffic density---while the road program followed its own
engineering logic. Yet the combined shock reallocated market access on
a national scale: 91~percent of districts gained market access between
1960 and 1986, because the road expansion's gains outweighed the rail
contraction's losses nearly everywhere
(Figure~\ref{fig:ma_choropleth}). A policy debate framed as ``closing
unprofitable lines'' was, in its spatial incidence, a large
redistribution of access to markets, with winners and losers that the
fiscal criteria never examined. This is the sense in which fiscal
consolidation acted as regional policy without being one, and it
suggests that the spatial incidence of infrastructure retrenchment
deserves the same scrutiny routinely applied to infrastructure
expansion.

\subsection{Market-Access Elasticities in Comparative Perspective}
\label{sec:comparative}

Our full-sample population elasticity of 0.046 (0.033) is an order of
magnitude below the benchmarks in the literature:
\citet{donaldsonhornbeck2016} report agricultural land-value
elasticities consistent with population responses near 0.25 in
nineteenth-century America, \citet{ahlfeldtfeddersen2018} find
approximately 0.35 for German high-speed rail, and \citet{gibbons2024}
approximately 0.3. Three measurement explanations for the gap were
examined and ruled out in the replication archive.\footnote{The
diagnostics are reproduced by the archive's reference-point,
spatial-standard-error, and elasticity-sweep programs; output files
accompany the package.} Re-anchoring market access at interior points
or at each district's largest settlement does not raise the estimate;
spatial standard errors leave the headline inference unchanged; and a
transshipment-cost bracket reshapes the market-access distribution
without changing the elasticity's order of magnitude.

What does move the level is the trade elasticity $\theta$: the
population coefficient ranges from roughly 0.37 at $\theta = 1$ to
0.01 at $\theta = 12$, and our main value $\theta = 4.55$ sits in
the interior of that range. The comparison with the literature is
therefore partly a comparison of calibrations, and we treat the
level of the population elasticity as provisional on that choice.
Two findings are robust to it. First, the sectoral contrast
documented below holds at every value of $\theta$ we examined.
Second, the context differs from the benchmarks in a substantive
way: the benchmark studies measure access expansions into settling
or growing economies, while we measure a reallocation within an
already-settled population, where mobility costs and accumulated
local capital plausibly damp population responses.

\subsection{Sector--Mode Specialization}
\label{sec:sector_mode}

The clearest elasticities in this paper are sectoral, not aggregate:
manufacturing production value responds at 0.259 (0.119) and
manufacturing wage mass at 0.330 (0.126), while agricultural outcomes
show no detectable response (Table~\ref{tab:sectoral_iv}). The cost
schedule of \citet{baumgartnerpalazzo1969} offers a candidate
explanation, stated in Section~2.4 as a hypothesis: rail is the
cheaper mode above roughly 500 to 1{,}000 tons per day of cargo
density and road below it (Figure~\ref{fig:cost_schedule}). A
modal shift from rail to road therefore lowered transport costs for
the low-density, higher-value shipments characteristic of
manufacturing, while raising them for the bulk flows characteristic
of agriculture. The counterfactual decomposition in Section~6 is
consistent with this reading but cannot sharpen it, because the
sectoral outcomes are not available at the frequency needed to
separate modal channels. Sector-specific market access measures,
with mode-specific cost calibrations, are the natural next step and
are left for future work.

\subsection{Limitations}
\label{sec:limitations}

Four limitations qualify the results. First, the pre-trends placebo
is not a clean null: post-reform market-access changes predict
1947--1960 population growth on the 235-district placebo subsample
(Table~\ref{tab:pre_trends}), and spatial standard errors sharpen
rather than soften that rejection. The placebo subsample differs
from the full sample, so selection and pre-trend explanations cannot
be separated with available data; readers should treat the aggregate
population estimates with corresponding caution. Second, the recent
in-migration coefficient in Table~\ref{tab:other_outcomes_iv} is
negative; both a displacement reading and a measurement reading are
consistent with it, and we do not adjudicate.
 Third, the
hypothetical-road instrument is weak for total market access once
road connectors are re-costed, so identification of the road
component leans on functional form more than we would like; the rail
component, by contrast, is strongly identified throughout. Fourth,
market access is a partial-equilibrium object: it holds destination
populations fixed and abstracts from price responses, so the
estimates are reduced-form gradients, not general-equilibrium
counterfactuals.

\section{Conclusion}
\label{sec:conclusion}

Between 1960 and 1991 Argentina closed roughly ten thousand
kilometers of railway and paved roughly eighteen thousand kilometers
of road, a rare episode in which a country simultaneously contracted
one national transport network and expanded another. We measure the
resulting change in every district's access to markets with a single
multimodal index, instrument it with the Larkin Plan's technical
study designations and with hypothetical road networks, and estimate
the response of population and economic activity across 312
districts.

Three findings summarize the evidence. Aggregate population
reallocation was modest: the full-sample elasticity is small,
sensitive to the trade-elasticity calibration, and not statistically
distinguishable from zero. Economic specialization was not:
manufacturing production and wages responded strongly where market
access rose, while agriculture did not respond, a contrast robust to
every calibration we examined and consistent with the modal cost
structure that made roads cheap for low-density cargo. And the
channel decompositions attribute the population response in
similar parts to rail and road shocks---with the rail component far
better identified---and roughly half of it to infrastructure located
in the district itself.

For policy, the episode illustrates that decisions taken on fiscal
grounds can carry first-order spatial consequences: the geography of
market access was redrawn without any actor being charged with
evaluating that redistribution. For research, the results caution
against transferring market-access elasticities across contexts
without attention to the trade-elasticity calibration and to whether
the underlying shock expands access into growing economies or
reallocates it within settled ones.

Three extensions follow directly. A population-weighted urban anchor
built from geocoded census records would settle the reference-point
question definitively. Sector-specific market access, calibrated to
mode-specific costs, would turn the sector--mode hypothesis into a
testable specification. And richer within-district data---stations
and depots rather than kilometers---would sharpen the local
decomposition that Section~7 begins.
